\section{Experimental Setup}
\label{sec:setup}

\subsection{Pre-registration}
\label{sec:prereg}

We registered hypotheses, effect sizes, arm definitions, decision rules, and per-role predictions
before data collection. Two amendments were applied predating the data they're related to.
Specifically, amendment A2 claimed that improvement on the manipulation check was our primary
outcome, rather than the benchmark gain. A3 acknowledges a pilot which led to the analysis of the
47-item test split and pre-registered Benjamini--Hochberg control for the six role checks.

The registered decision rule is conjunctive. A role passes only if its paired
operation-versus-neutral contrast survives FDR control at $q = 0.05$ across the six-role family
\emph{and} reaches $d \ge 0.40$; a role is recorded as moving a benchmark only above a
leave-one-out contrast of 0.05. All intervals are 95\% paired bootstrap intervals over 2,000
item-level resamples. The scoped study is underpowered for the original council-gain hypothesis by
design --- at $n = 1{,}000$ the detectable effect is $\approx 0.048$ against a SESOI of 0.03 --- so
that hypothesis is reported as \textbf{not tested}, never as a null. The manipulation check is
differently placed: at $n = 200$ paired, the minimum detectable effect at 80\% power is
$\approx 0.20$, so the registered $d \ge 0.40$ threshold sits comfortably above the detection floor.

\subsection{Arms and cost matching}
\label{sec:arms}

The registered design has four arms: a full-operation council, a persona-label-only council, a
length-matched neutral council, and a cost-matched single-agent baseline. The council arms differ
only in the prompt (\S\ref{sec:variants}); model, decoding, output contract, and aggregation are
identical.

\textbf{The fourth arm was collected, but not as registered.} The registered baseline was
cost-matched. So, what we ran is the single-agent arm on the same model the councils use with one
direct call, making it neither cost-matched nor call-matched to the councils. The preregistered
cost-matched contrast remains uncollected. A council deficit against this baseline is therefore a
deficit against something that also costs less: a stronger practical claim and a weaker answer to
the registered question. \S\ref{sec:attribution} states the consequence for the attribution
baseline.

Two benchmarks are collected, both materialised by script from pinned immutable revisions and
verified by content hash so the corpus is reproducible byte-for-byte: HaluEval's QA split
\citep{li2023halueval} at \texttt{207f479c}, expanded to 20,000 binary items, and TruthfulQA's
generation split \citep{lin2022truthfulqa} at \texttt{741b8276}, 817 items. We use both for
non-commercial research evaluation, following the terms stated in their licenses. No personal data
is contained and we redistribute with the original license with attribution, with license
identifiers and URLs recorded per benchmark in \texttt{data/materialize\_manifest.json}.

\textbf{Accuracy analysis is limited to HaluEval.} Because the TruthfulQA generation split does not
provide an answer index (grading uses a judge model), the gold answer index is necessarily undefined
for all TruthfulQA questions. We classified TruthfulQA as free-form in preprocessing but did not
recognise before collection that the registered accuracy analysis needs a gold answer index, which
the generation split lacks. Examples from the TruthfulQA benchmark are, however, reflected in
behavioural metrics shown in \S\ref{sec:results} (e.g., refusal, abstention rates). All accuracy
metrics are, in turn, limited to HaluEval and reported as such (total examples), everywhere it
appears. For the benchmark arms, we use 1,000 HaluEval examples per arm (see
Appendix~\ref{app:splits}).

\subsection{Models}
\label{sec:models}

For our benchmark arms, we host \texttt{gpt-5.6-terra} in Microsoft Foundry (eastus2 Global
Standard), and for cross family replication in \S\ref{sec:replication} we host \texttt{grok-4.3} in
the same. Our pre-registered cost-matched arm uses \texttt{gpt-5.6-luna}, but as we switch Luna in
the collected single agent arm to Terra (\S\ref{sec:arms}), results do not include Luna. For all
arms and roles, we use \texttt{temperature=0}, 1,024 output tokens, \texttt{seed=1337},
\texttt{timeout=120s}. Of all provenance facts, only the following two are asserted (not inferred)
--- and only these two: (i) We specify the deployment name as the ID in our records. This deployment
is known to have used model version \texttt{2026-07-09}. The deployment name does not reference this
version: we thus specify the version in the manifest and do not attempt to infer it from the records
themselves; (ii) We specify the serving route (vendor, region, deployment mode), not infer it, since
our collection script accepts a provider argument. All other provenance facts are inferred (prompts
via hash, simulated/reality from provider attestation, splits by content).

\paragraph{Compute.}
All model calls were hosted as API calls, without using any GPU cluster nor hardware on the
provider's side. In collection, we issued 35,700 model calls. This includes three council arms at
1,500 items each $\times$ 6 different roles, another single agent arm at 1,500 items, and then three
prompt conditions at 200 items $\times$ 6 roles in each of \texttt{gpt-5.6-terra} and
\texttt{grok-4.3}. 35,208 (2.2\% absent rate) returned usable records, with recorded usage totalling
about 8.09M output tokens (\$128.66). Each arm in the council costs six calls per item: around 3,350
input and 1,428--1,659 output tokens per item (\$0.0239--\$0.0267), with mean per-agent latency at
3.1--4.5\,s; the single-agent baseline averages 2.44\,s (median 2.23\,s) and \$0.001792 per call.
Per-run cost, token, and latency breakdowns are in each run's \texttt{evaluation.json} under
\texttt{cost}. The synthesizer is the only trained component: 81--108 optimizer steps, 34.7--55.3\,s
wall clock per run on one Apple Silicon laptop (MPS, no CUDA). Every analysis in this paper re-runs
from the shipped records in under a minute on the same machine, with no API access.

\subsection{Attribution estimator}
\label{sec:attribution}

\textbf{In this work, we specify the value of the empty coalition.} As specified in our pre-reg, we
define $v(\emptyset)$ as the direct answer base model evaluated on identical per-sample inputs, with
records provided by the single-agent arms and joined into our council records. Note that we are able
to join 1,134/1,500 records per arm (including all 1,000 scorable records): we do not infer missing
values for 366 TruthfulQA records where the baseline abstained (or returned free-form text with no
answer index), we simply report these as missing. (Either way, leave-one-out does not consider the
empty coalition at all.)

For per-role attribution we perform leave-one-out by retraining on each coalition (seven trainings
total rather than the 64 required for a full Shapley computation, which is computationally
prohibitive at this scale; using a masking approximation to improve tractability deviates from
retraining by more than the effect size of interest, see Appendix~\ref{app:attribution-failures}).
Explicitly, leave-one-out is not a Shapley value and we do not claim it to be: in particular, it is
not efficient, so per-role contributions do not have to add up to $v(A) - v(\emptyset)$, and it
attributes all role--council interactions to each role. Writing $v$ in Harsanyi dividend form
$v(S) = \sum_{\emptyset \ne T \subseteq S} d(T)$, leave-one-out evaluates to
$v(A) - v(A \setminus \{i\}) = \sum_{T \ni i} d(T)$ while the Shapley value is
$\sum_{T \ni i} d(T)/|T|$, so the two differ by
$\sum_{T \ni i,\, |T| \ge 2} \frac{|T|-1}{|T|} d(T)$: the standalone term $d(\{i\})$ cancels and
leave-one-out charges role $i$ the whole of each interaction rather than its $1/|T|$ share.
Substitutable roles are negative dividends, so under the pairwise substitutability registered in
advance this difference is negative and every leave-one-out contrast is biased downward relative to
the Shapley value. The sign is imported rather than measured --- the seven coalitions evaluated here
do not identify $d(T)$ for $|T| \ge 2$ (Appendix~\ref{app:third-finding}) --- but the bias runs
against a positive verdict either way, so a role exceeding the 0.05 leave-one-out (LOO) threshold
has passed a conservative test.

Provenance is derived rather than asserted. Refusals, truncations, malformed responses, and retries
are counted and published per role rather than imputed. In the pilot and manipulation-check runs
reported here, those rates are zero across 1,128 pilot calls and the 3,600 calls; the benchmark arms
refuse only on TruthfulQA items, entering no accuracy contrast (Appendix~\ref{app:refusal}).
Appendix~\ref{app:provenance} gives the mechanism.
